\section{静的な多体力}
\label{sec:force}

地位は\dstrec{}である。この章の関数は時刻を受け取らない。加速度も返さない。

\subsection{重力だけの動径項}

\(\kappa\) は入力 \(G m_1 m_2\) である。
\[
\texttt{gravityLayerTerm}(\kappa,\ell,r)
=-\frac{\kappa}{\gamma_s^{\mathrm{eq}}(\ell/(2r))\,r^2}.
\]
正の値は分離を増やす。\(\kappa>0\)、\(r>0\)、\(\gamma_s^{\mathrm{eq}}<0\) ならば項は正である。同じ仮定で \(\gamma_s^{\mathrm{eq}}>0\) ならば項は負であり、分離を減らす向きである。

点 \(p\) が点 \(c\) から受けるベクトルは
\[
\texttt{gravityLayerPush}(\kappa,\ell,c,p)
=\frac{\texttt{gravityLayerTerm}(\kappa,\ell,\lVert p-c\rVert)}{\lVert p-c\rVert}\,(p-c)
\]
である。外向き単位分離との内積が、動径項に一致する。反対称性
\[
\texttt{gravityLayerPush}(\kappa,\ell,c,p)
=-\texttt{gravityLayerPush}(\kappa,\ell,p,c)
\]
を、二体テストの等式として使う。

有限集合 \(S\) の粒子 \(i\) に働く力は
\[
\texttt{gravityLayerForce}_i
=\sum_{j\in S\setminus\{i\}}
\texttt{gravityLayerPush}(G m_i m_j,\ell,\mathrm{pos}_j,\mathrm{pos}_i)
\]
である。和は \texttt{erase i} であり、自己項は無い。実装手順は次の順にする。

\begin{enumerate}
\item 二体のスカラー \texttt{gravityLayerTerm} を、引力層と斥力層のサンプルで評価する。
\item 二体のベクトルが反対称であることを見る。
\item 三体で、粒子 \(i\) の和に \(i\) 自身が現れないことを見る。
\item \(N\) 体へ広げる。各項の失敗（\(r=0\) または \(\gamma=0\)）は、和を部分的に返さず、呼び出し全体を拒否する。
\end{enumerate}

\subsection{二チャネル}

識別子 \texttt{pairChannelForce} は分離を増やす向きのスカラーである。
\[
-\frac{\kappa}{\gamma_G r^2}+\frac{\mathrm{coul}}{\gamma_E r^2}.
\]
\(\kappa=G m_1 m_2\)、\(\mathrm{coul}=k q_1 q_2\) である。\(\gamma_G,\gamma_E,r\) が非零のとき、この量が零になることと
\[
\kappa\,\gamma_E=\mathrm{coul}\,\gamma_G
\]
は同値である。二つの干渉因子を固定したままでは、この条件に半径は現れない。正の半径がすべて零か、どれも零でないかである。\(\gamma\) を \(r\) から計算し直した合成関数の零を、平衡半径として探索するループは書かない。

重力項は、電荷項を \(0\) にし \(\gamma_E=1\) とした二チャネル力と一致する。電磁項を \texttt{gravityLayerForce} の和へ混ぜない。二チャネルが要る呼び出しは、別関数 \texttt{pair\_channel\_force} を使う。

\(\kappa>0\) かつ \(\mathrm{coul}<0\) の異種対では、干渉因子が同符号なら分離を増やす力は零にならない。両方の因子が正なら力は負であり、両方とも負なら力は正である。実装はそのサンプルをテストに残す。スペクトルも束縛エネルギーも、この関数の戻り値にしない。

\subsection{クーロン項と点電荷}

スカラー
\[
\texttt{coulombForce}(k,e,\gamma_s,r)=\frac{k e^2}{\gamma_s r^2}
\]
の符号は \(\gamma_s\) の符号である。円軌道の読み \(v^2=ke^2/(m\gamma_s r)\) は、このスカラーと \(r/m\) の積として書かれている。その \(m\) は呼び出し側のパラメータであり、剛性からも \(M\) からも計算しない。

金属側の点電荷力は、\(\ell=1\) に正規化した別モジュール \texttt{ElectronMetal} である。\texttt{pullRate} は \(ke^2/\ell^3\) を単位に含んだ量であり、\texttt{pointForce} は \(k\) と \(e\) を別引数に取らない。
\begin{align}
\texttt{vecPhase}(w)&=\frac{1}{2\lVert w\rVert},\\
\texttt{pullRate}(x)&=\frac{8x^3}{\gamma_s^{\mathrm{eq}}(x)},\\
\texttt{pointForce}(Q,c,p)
&=-\bigl(Q\cdot\texttt{pullRate}(\texttt{vecPhase}(p-c))\bigr)\,(p-c).
\end{align}
\(r=\lVert p-c\rVert>0\) かつ \(\gamma_s^{\mathrm{eq}}>0\) かつ \(Q>0\) のとき、\(p\) に働くベクトルは \(c\) の向きを向く。この正規化での大きさは \(|Q|/( \gamma_s^{\mathrm{eq}} r^2)\) に一致する。確認は \(r=1/(2x)\) と \(8x^3\cdot r=4x^2\) と \(4x^2=1/r^2\) を使えばよい。\(r=0\) または \(\gamma_s^{\mathrm{eq}}=0\) の呼び出しは拒否する。\(\ell\neq 1\) の計算で \texttt{vecPhase} を流用しない。長さを戻すなら、位相を \(\ell/(2r)\) に取り、前置因子を単位に合わせて別関数にする。その関数は記録の \texttt{pointForce} と別名にする。

結晶の \texttt{siteForce} は \texttt{ElectronCrystal} にあり、有限集合 \(T\) 上のベクトル和
\[
\sum_{t\in T}(-q)\,\texttt{pointForce}(Q_t,\,p-\texttt{latticeVec}(b,h,t),\,x)
\]
である。\((-q)\) はベクトルへのスカラー倍であり、内積ではない。\(q\) は試験電荷、\(b\) は正規直交基底、\(h\) は格子間隔である。時間は進まない。無限格子和、遮蔽、凝集エネルギーは、この和の定義に入っていない。

\subsection{チャネルを一つの和にしない}

重力の \(N\) 体和と、電磁の \(N\) 体和は、別の関数として残す。呼び出し側が両方を欲しがるときは、二つのベクトルを受け取ったあとで足す。足し算そのものは第\ref{sec:contract}章の作業段に属する。Lean の \texttt{gravityLayerForce} は電磁項を含まず、\texttt{pointForce} の和は重力項を含まない。
